\section{Related Work}

\textbf{Terrain representations.}
Occupancy maps such as OctoMap \cite{hornung2013octomap} and signed-distance fields such as
Voxblox \cite{oleynikova2017voxblox} provide geometrically meaningful state for planning.
Elevation maps are especially efficient for legged motion and have enabled robust
rough-terrain locomotion \cite{fankhauser2018rough}. These representations generally encode
sensor noise but do not expose the decision-specific probability that a proposed contact
will fail after occlusion. \method{} retains the efficient $2.5$D interface while learning
how map error evolves when evidence becomes stale.

\textbf{Learning for locomotion.}
Massively parallel simulation can produce capable locomotion policies in minutes
\cite{rudin2022minutes}; perceptive policies have since demonstrated impressive outdoor
mobility \cite{miki2022wild}. Learned policies can internalize some observation corruption,
but their confidence is difficult to inspect or constrain. Our design instead learns a
mapping module and gives an optimization-based controller an interpretable chance-risk
budget.

\textbf{Uncertainty-aware planning.}
Aleatoric and epistemic uncertainty capture different failure sources
\cite{kendall2017uncertainty}. Informative planners have used uncertainty to select sensing
actions \cite{schmid2020approach}, while gaze-control methods reason about what the robot
should observe \cite{wellhausen2020terrain}. We focus on the complementary interval after
an observation is lost: memory preserves the last evidence, calibration tells the planner
how quickly to stop trusting it, and contacts selectively repair the estimate.
